\section{Stochastic satisfiability: ordered evaluation and inference}
\label{sec:p15-ssat-annex}

These proofs establish the declared finite ordered calculus, biased
independent extension and explicit inference/sampling constructions.
The directly read source is the January 2000 author manuscript; these
proofs identify their fair source and their independently declared weighted scope.

\begin{annexitem}{annex:p15.ssat-native-correspondence}{thm:p15-ssat-native-correspondence}{The actual ordered-map and prefix comparison}
\begin{proof}
An order-preserving bijection between finite variable lists carries each
assignment to the assignment with the same Boolean values at the renamed
variables. Induction on literals and clause truth shows that this bijection
preserves the matrix payoff. Restriction commutes with it: true clauses
are deleted in both matrices and false literals removed in both. The
source's renumbering uses precisely this bijection on the remaining list,
and transports the quantifier map with it; it does not permute positions.

Induct on the number of remaining variables. At either terminal, both
definitions return the same zero or one, independently of unused variables.
Otherwise the order-preserving correspondence gives matching outermost
quantifiers and corresponding two cofactors. The induction hypothesis
identifies each child value. The native maximum and minimum are exactly
the prefix maximum and minimum, and its arithmetic mean is the weighted
rule at parameter \(1/2\). Applying those operators gives equality of the
parent values. This proves both renaming invariance and the claimed native
map/prefix equality for every finite input. Comparing equal values with
the same threshold proves the decision correspondence. No step identifies
different operator orders or changes a nonfair weight to \(1/2\).
\end{proof}
\end{annexitem}

\begin{annexitem}{annex:p15.ssat-strategy-semantics}{thm:p15-ssat-strategy-semantics}{Recursive soundness, completeness and nonanticipating strategies}
\begin{proof}
Restriction preserves the truth of a CNF under the assigned value:
clauses with a true literal are already satisfied; false literals
contribute nothing to their clause; a resulting empty clause is false.
An empty matrix is true under all continuations, and an empty clause
is false under all continuations. Thus the two terminal rules agree
with the matrix payoff for every remaining prefix.

Induct on the remaining prefix length. At an existential node there
are two available actions. The controller chooses one before seeing
any later draw, so the largest achievable continuation success is
the maximum of the two already determined child values. Choosing an
argmax and the inductively constructed continuation strategies attains
it. At a universal node the opposing player chooses an argmin, giving
the minimum. At a random node the fixed independent outcome weights
are \(1-p,p\), so finite total expectation gives their weighted
sum. After the random value is revealed the continuation strategies
may differ in the two branches, which is exactly the legal information
condition. No earlier decision is permitted to use that branch bit.

More precisely, backward induction constructs a maximizing strategy
which guarantees at least the node value against every minimizing
continuation, and a minimizing strategy which guarantees at most it
against every maximizing continuation. At max nodes the maximizing
player selects the largest child; the minimizing player installs its
guaranteeing continuations in both children. At min nodes reverse
these roles. At chance nodes install both children's strategies and
average their guaranteed payoffs. This proves matching guarantees
and attains the value for the constructed strategy pair. A mixed
controlled action is a convex combination of the two child payoffs,
so cannot improve on their maximum or minimum. Pure backward actions
therefore suffice.

When there are only existential and random nodes, earlier existential
values are determined recursively by their own tables. Eliminating
them as independent table arguments leaves exactly the preceding
random-history arguments claimed in the body. With universal nodes,
all permitted earlier opposing/chance observations remain available.
The construction fills the entire finite history tree, including null
histories; no division by a history probability occurs. For \(p=0\)
or \(p=1\), a child of weight zero simply does not affect its parent.

Max, min and normalized nonnegative weighted sums preserve \([0,1]\).
Every nonterminal application removes one quantifier and branches into
at most two smaller problems. Thus the complete tree is finite,
including unused variables whose two cofactors agree. At an empty
prefix, all variables have been restricted, so the CNF is terminal.
The induction proves existence of a derivation computing the value
and that every such derivation computes that same value. This is
soundness and completeness of the finite evaluation rules.
\end{proof}
\end{annexitem}

\begin{annexitem}{annex:p15.ssat-prefix-order}{thm:p15-ssat-prefix-order}{Like operators commute; choice and observation do not}
\begin{proof}
Let \(f(a,b)\) be the value of the common suffix after fixing two
adjacent variables. Two maxima give \(\max_{a,b}f(a,b)\) in either
order; two minima give \(\min_{a,b}f(a,b)\). Two independent random
quantifiers give the finite sum
\(\sum_{a,b}\mu_a\nu_bf(a,b)\); distributivity and interchange
of the two finite sums prove equality. The biases belong to their own
variables and are not exchanged. Applying the unchanged earlier
prefix to these pointwise equal continuation values preserves equality.

For \(x\leftrightarrow y\), a controlled \(x\) before the
random draw must be constant. Choosing one succeeds with probability
\(p\), choosing zero with probability \(1-p\); the optimum is
their maximum. After \(y\) is observed the policy \(x=y\)
is legal and succeeds with probability one. At \(p=1/4\) these
give \(3/4,1\), separated by \(4/5\).

For every fixed action \(a\), \(f(a,b)\leq\max_{a'}f(a',b)\)
for every \(b\). Multiply by nonnegative weights and sum, then
maximize the left side over \(a\). This proves the general information-
relaxation inequality.

For completeness, restrict the source's three-clause example. At
\(x=0\), its matrix reduces to \(y\land(y\lor\bar z)\),
whose best \(y\) is one, giving value one. At \(x=1\) it
reduces to \((y\lor\bar z)\land\bar y\), giving value
\(1/2\) when \(y=0\). Averaging gives \(3/4\).
For \(\ssatrandom x\ssatrandom z\exists y\), the assignments
\((x,z)=(0,0),(0,1),(1,0)\) have satisfying extensions and
\((1,1)\) does not, again giving \(3/4\). With \(y\)
outermost, \(y=1\) succeeds precisely when \(x=0\), giving
\(1/2\), while \(y=0\) requires \(x=1,z=0\), giving
\(1/4\). The maximum is \(1/2\). Every displayed calculation
therefore respects its own prefix.
\end{proof}
\end{annexitem}

\begin{annexitem}{annex:p15.ssat-unit-pure}{thm:p15-ssat-unit-pure}{Structural simplifications without illegal information}
\begin{proof}
Each value operator is monotone in both children and positively
homogeneous: for \(s\geq0\), combining \(sa,sb\) equals
\(s\) times the combination of \(a,b\). The max and min identities
follow from order preservation under nonnegative scaling; the random
identity is distributivity. By induction, these facts hold for every
unchanged prefix applied to its leaf continuation values.

Suppose a unit forces \(x=b\). For every fixed history before
\(x\), the opposite cofactor contains an empty clause and has
value zero. At \(x\), existential max therefore gives the good
cofactor value, universal min gives zero, and random averaging gives
\(p_b\) times the good cofactor. The random factor is the same
at every earlier history by the fixed-independent-parameter hypothesis.
Positive homogeneity pulls it through every earlier operator, proving
the displayed identity even for an inner unit variable. If \(p_b=0\),
the value is identically zero; no inverse weight is needed.

For a pure variable with satisfying sign \(b\), pointwise matrix
truth at \(x=1-b\) implies matrix truth at \(x=b\): changing
it cannot falsify any clause. Applying the unchanged suffix's monotone
operators preserves this dominance for every earlier history. Thus
existential max at \(x\) selects \(b\), while universal min
selects \(1-b\). Applying the earlier prefix preserves these
equalities. This is not an unrestricted reordering argument: it uses
a dominant cofactor at every history. Random averaging instead retains
both contributions, and \(\ssatrandom^{1/4}x\;(x)\) explicitly
gives \(1/4\), disproving random purification.

A tautological clause is true under every total valuation, so deleting
it leaves the Boolean payoff unchanged; induction through the prefix
preserves its value. An unused variable has equal cofactors of value
\(v\); all three operators return \(v\), including deterministic
random endpoints. This proves the remaining simplifications.
\end{proof}
\end{annexitem}

\begin{annexitem}{annex:p15.ssat-window-correctness}{thm:p15-ssat-window-correctness}{Complete threshold-window induction with endpoints}
\begin{proof}
Use induction on the number of prefix variables. Terminals return
their exact value, satisfying any window. The unit/pure identities
already proved justify their recursive smaller calls. A positive
random-unit factor rescales every comparison and the exactness interval
by the same \(s>0\), so the transformed contract gives the root
contract. A zero factor returns exact zero. All other transformations
leave the contract unchanged.

Let true child values be \(a,b\), and returned first value \(u_0\).
Under the induction contract, \(u_0<L\) holds exactly when
\(a<L\), \(u_0>H\) exactly when \(a>H\), and a return
in the closed window equals \(a\).

At an existential node, if \(u_0>H\), then \(a>H\) and
\(\max(a,b)>H\), so the early return is valid. Otherwise
\(a\leq H\), and either \(u_0<L,a<L\), or \(u_0=a\)
in the window. Put \(L'=\max(L,u_0)\leq H\). If
\(b<L'\), the second return is below \(L'\). When
\(u_0\geq L\), \(L'=u_0=a\), so the returned maximum
is exactly \(a=\max(a,b)\). When \(u_0<L\), both true
children and both returned children are below \(L\), so their
maximum is correctly below \(L\). If \(L'\leq b\leq H\),
the second return is exactly \(b\), and \(b\geq a\) by
the same two first-child cases, so the maximum is exact. If
\(b>H\), the second return is above \(H\), as is the true
maximum. This covers every case, including equality with both endpoints.

At a universal node, if \(u_0<L\), then \(a<L\), making
the true minimum and returned value below \(L\). Otherwise
\(a\geq L\), and either \(u_0=a\leq H\), or both
\(u_0,a>H\). Put \(H'=\min(H,u_0)\geq L\). If
\(b>H'\), then when \(u_0\leq H\), \(H'=u_0=a\)
and the returned minimum is exactly \(a\). When \(u_0>H\),
both true and returned children are above \(H\), so their minima
are above it. If \(L\leq b\leq H'\), the second return is
\(b\), which is no larger than \(a\), so the returned minimum
is exact. If \(b<L\), both the true and returned minima are
below \(L\). These are all universal cases.

At a random node with \(p=0\) or one, the positive-weight cofactor
is exactly the whole value, so its unchanged contract suffices. For
\(0<p<1\), put \(q=1-p\),
\(L_0=(L-p)/q\), \(H_0=H/q\). These satisfy \(L_0\leq H_0\).
If the first return satisfies \(qu_0+p<L\), it is below
\(L_0\), so \(a<L_0\); since \(b\leq1\),
\(qa+pb\leq qa+p<L\). The returned \(qu_0\) is also
below \(L\). If \(qu_0>H\), the first return is above
\(H_0\), so \(a>H_0\); since \(b\geq0\), the true
value is above \(H\), and the returned \(qu_0\) is above
it as required. If neither test holds, then
\(L_0\leq u_0\leq H_0\), forcing \(u_0=a\) by the
induction contract. The second window
\(L_1=(L-qa)/p\), \(H_1=(H-qa)/p\) is well ordered.
If \(b\) is below, within, or above it, the return is respectively
below, exact within, or above it. Multiplication by positive \(p\)
and addition of \(qa\) give precisely the three root-window
cases. No null-weight division or unexamined equality case is used.

Each call removes a variable, so induction also proves termination.
Every true value lies in \([0,1]\), making the \((0,1)\) call
exact. At \((\theta,\theta)\), a value below the threshold returns
below it, equality returns exactly it, and a value above returns above
it. Thus comparison by \(\geq\) is exactly the intended decision.

For the printed-source countercase, the previous calculation gives
first child \(1/4\) and second child \(1/2\) under outer
\(\exists y\). A printed \(u_0\geq H\) early return at
\(H=1/4\) returns the endpoint instead of a value above it,
although the actual root is \(1/2\). The corrected strict test
does not return there. This establishes the actual failing boundary.
\end{proof}
\end{annexitem}

\begin{annexitem}{annex:p15.ssat-certificate-complexity}{thm:p15-ssat-certificate-complexity}{Finite exact certificates, rational space and QBF reduction}
\begin{proof}
The certificate checker verifies each recorded restriction and prefix
step, then the appropriate arithmetic equation. Induction from its
terminal leaves shows that every accepted node's value is the genuine
continuation value. Hence an accepted root is exact. Conversely build
the canonical full restriction tree, labeling every node by the
finite evaluation recurrence. It satisfies every check. Its depth is
at most \(n\), with at most \(2^{n+1}-1\) nodes. Thus complete
finite certificates exist, without a polynomial-size claim.

Finite literal-set membership/restriction tests are polynomial in
the explicitly supplied formulas. Rational equalities and comparisons
can be checked by integer cross multiplication; schoolbook arithmetic
is polynomial in the explicit operand bit lengths. All recorded
operand bits count toward certificate length. Therefore checking is
polynomial in that length plus the input, not necessarily in the input
alone. This is the mathematical checker guarantee, not a claim that
a parser or executable checker has been tested.

Write the independent random probabilities as \(a_i/b_i\),
with \(0\leq a_i\leq b_i\), \(b_i\geq1\). For every suffix
use common denominator \(D\) equal to the product of denominators
of its remaining random variables. By induction its value has numerator
\(N\) with \(0\leq N\leq D\). Max and min choose between
numerators over the same suffix denominator. A random step combines
them as
\[
 N'=(b_i-a_i)N_0+a_iN_1,
 \qquad D'=b_iD.
\]
The bit length of every such product denominator is at most the sum
of input denominator bit lengths, plus a constant. Numerators have
the same bound. Depth-first evaluation stores at most \(n\) stack
frames, restricted formulas of polynomial size and these polynomial-
length integers. Its space is polynomial. The encoded threshold is
compared by one further polynomial-length cross product.

For hardness start with any encoded closed QBF. If its Boolean matrix
is not CNF, introduce a fresh output variable per Boolean gate and
the full equivalence clauses described in the projection/inference
proofs. Require the output true and append all gate variables to an
innermost existential block. For each fixed original valuation there
is exactly one consistent gate assignment, and an extension satisfying
the output unit exists exactly when the original matrix is true.
Thus every original QBF prefix has the same truth value after this
linear-size CNF transformation.

Now replace each universal quantifier by a fair random quantifier and
retain the existentials. The resulting SSAT uses threshold one.
Induct from its leaves: a max of two values in \([0,1]\) is one
exactly when at least one is one; their fair mean is one exactly when
both are one. These are precisely existential and universal truth
conditions in the original QBF, including constant terminals and
unused variables. Hence the original QBF is true exactly when this
SSAT instance is positive. This is a polynomial many-one reduction.
The already proved TQBF \(\mathsf{PSPACE}\)-completeness in the
book, together with membership above, establishes the stated result.
\end{proof}
\end{annexitem}

\begin{annexitem}{annex:p15.ssat-projections}{thm:p15-ssat-projections}{Counting, maximizing a fixed choice and contingent projection}
\begin{proof}
Repeated expansion of random averages gives a sum over every random
assignment of its product weight times the terminal truth indicator.
This proves the random-only case. Fair weights make every complete
assignment's weight \(2^{-n}\), giving the model-count ratio.
An initial existential block applies a maximum over one fixed assignment
to that block before the remaining weighted sum, giving the second
formula. A final existential block has value one exactly for random
assignments admitting a satisfying extension and zero otherwise;
averaging those indicators gives the third formula.

For final existential auxiliaries, only the event that at least one
extension exists matters. Equality of those projected events on every
random assignment proves probability preservation, regardless of the
number of extensions. In contrast, random auxiliaries average over
their extensions. A uniquely forced fair Boolean auxiliary has exactly
one of its two outcomes compatible with its gate equivalence, so gives
factor \(1/2\); \(g\) such independent fair auxiliaries give
factor \(2^{-g}\). A genuinely unconstrained auxiliary contributes
the sum of its two weights, namely one, and so is a different case.

For explicit CNF gate projection, \(z\leftrightarrow a\land b\)
is encoded by \((\bar z\lor a),(\bar z\lor b),
(z\lor\bar a\lor\bar b)\). If \(z=1\), the first two force
both inputs true; if both inputs are true, the third forces \(z=1\).
The other input cases force \(z=0\), proving equivalence and uniqueness.
For \(z\leftrightarrow a\lor b\), use
\((z\lor\bar a),(z\lor\bar b),(\bar z\lor a\lor b)\):
a true input forces \(z=1\), and both false inputs force zero.
For \(z\leftrightarrow\bar a\), use
\((\bar z\lor\bar a),(z\lor a)\), forcing opposite values.
Topological gate evaluation proves unique extension for an arbitrary
finite circuit; an output unit then projects exactly its true-input
event. These clauses justify both the QBF transformation and the
event-circuit encoding, rather than assuming SAT equisatisfiability
alone preserves weighted counts.
\end{proof}
\end{annexitem}

\begin{annexitem}{annex:p15.ssat-bn-translation}{thm:p15-ssat-bn-translation}{Acyclic row selection and exact Bayesian joint events}
\begin{proof}
Fix any full coin assignment. At a root node the unique empty-parent
pattern selects its one coin, fixing \(A_i\). Induct in topological
order: all parents have already been fixed, exactly one parent pattern
is true, and the equivalence fixes \(A_i\) to the coin for that row.
Thus there is exactly one network-variable assignment. The gate
equivalences just proved similarly determine all auxiliaries. Every
coin assignment has this extension, including assignments of zero
product probability.

For a proposed network assignment \(a\), the event that this
extension is \(a\) imposes, for each \(i\), the single active
constraint \(C_{i,\operatorname{pa}_i(a)}=a_i\). It leaves every
other row coin free. Independence gives its probability
\[
 \prod_{i=1}^m
 \begin{cases}
 t_{i,\operatorname{pa}_i(a)},&a_i=1,\\
 1-t_{i,\operatorname{pa}_i(a)},&a_i=0.
 \end{cases}
\]
This is exactly the Bayesian-network factorization. Unused row coins
sum to one, even for deterministic zero/one row parameters. Therefore
the induced distribution equals the specified joint distribution.

For event \(B\), the final existential matrix succeeds exactly for
coin assignments whose unique extension satisfies \(B\). The
random-exist projection theorem makes its value the probability of that
event, proving the joint-event equality. Finite conditional probability
is the joint-event ratio only at positive \(P(E)\), giving the
displayed two-evaluation formula. If \(P(E)=0\), also
\(P(H\land E)=0\); their ratio is \(0/0\) and neither the
encoding nor a numerical convention supplies a posterior.

There is one coin per explicitly listed row. A row pattern has length
equal to that node's parent count. Combining the row expressions and
event circuits with linear-size gate equivalences is polynomial in
that explicit representation and its rational bits. No polynomial
claim is made relative to a compressed object hiding exponentially
many rows.

For the worked two-node network, the selected \(X=1\) row contributes
\((1/5)(4/5)=4/25\) to \(X\land Y\); the other row contributes
\((4/5)(1/10)=2/25\) to \(Y\), giving \(6/25\) total and
posterior \(2/3\). Treating the two marginal values independently
instead gives \(6/125\) for the joint event, and posterior \(1/5\),
which is the stated failure of that alternative encoding.
\end{proof}
\end{annexitem}


\begin{annexitem}{annex:p15.ssat-sampling-bound}{thm:p15-ssat-sampling-bound}{Fixed-policy concentration, uniform selection and sampled recursion}
\begin{proof}
For a fixed policy, each sampled full random assignment induces a
Bernoulli success indicator \(Z_j\) of mean \(v_s\), and the
indicators are independent. Thus the empirical average is unbiased.
Here is a full concentration argument. For a Bernoulli variable \(Z\)
of mean \(v\), set \(F(t)=\log E[e^{t(Z-v)}]\).
The expectation is positive for every real \(t\), including \(v=0,1\).
Differentiating the two-term finite sum gives \(F(0)=F'(0)=0\)
and \(F''(t)=a_t(1-a_t)\leq1/4\), where \(a_t\in[0,1]\)
is the exponentially tilted success probability. Taylor's integral
formula, for either sign of \(t\), gives \(F(t)\leq t^2/8\).

Independence yields
\(E[e^{t\sum_j(Z_j-v)}]\leq e^{wt^2/8}\).
For \(t>0\), Markov's inequality applied to that nonnegative
exponential gives
\[
 P\!\left(\sum_j(Z_j-v)>w\varepsilon\right)
 \leq \exp(-tw\varepsilon+wt^2/8).
\]
Choosing \(t=4\varepsilon\) gives \(e^{-2w\varepsilon^2}\).
Using \(t=-4\varepsilon\) for the lower tail gives the same bound.
Adding the two tail probabilities proves the fixed-policy inequality.
Union-bounding over the \(P\) complete tables gives the simultaneous
inequality; independence between different policies is not required.

On the uniform good event, let \(s_*\) maximize true value and
\(\widehat s\) maximize sample value. Then
\[
 V-\varepsilon\leq\widehat v_{s_*}
   \leq\widehat v_*\leq V+\varepsilon,
 \qquad
 v_{\widehat s}\geq\widehat v_* -\varepsilon\geq V-2\varepsilon.
\]
Solving \(2Pe^{-2w\varepsilon^2}\leq\delta\) proves the stated
sufficient sample size. If a returned local-search policy is at most
\(g\) below the sampled optimum, the same inequalities give true
value at least \(V-g-2\varepsilon\). No bound on \(g\) has
been inferred from the act of hillclimbing.

For sampled recursion, a nonempty sample multiset at an outer random
node splits into multisets of sizes \(w_0,w_1\). Its empirical
average is the sum of the two averages weighted by \(w_b/w\).
If \(w_b=0\), its contribution is zero, and its undefined empty
average is never computed. Future policy entries in two different
random-history branches are distinct permitted table entries, so the
two child empirical optima can be combined. At an existential node
one decision is shared by all sample histories reaching that node;
maximizing its two cofactor values is precisely the legal empirical
optimization. Terminal sample payoffs are zero or one. Induction on
the remaining prefix proves equality with the maximum empirical
value of complete legal policies. Any unused table entries can be
completed without affecting the sample. The prefix strictly decreases,
proving termination. This procedure respects the empirical joint
history distribution, not a product of fitted empirical marginals.
\end{proof}
\end{annexitem}

\begin{annexitem}{annex:p15.ssat-sampling-overfit}{thm:p15-ssat-sampling-overfit}{The majority-per-history construction and necessary exponential sample size}
\begin{proof}
After the \(r\) fair initial bits are observed, the controlled
\(x\) still precedes the independent fair final \(z\). Either
controlled value matches it with probability \(1/2\), so every
complete policy has true value \(1/2\). There are \(2^r\)
histories, each permitting one independent choice of a table bit,
hence \(2^{2^r}\) policies.

On the sample, the table entry for history \(h\) chooses the
majority of its observed final bits. Its number of successes in that
group is \(\max(N_{h,0},N_{h,1})\), proving the normalized sum
formula. Fix a sampled index \(j\). Each other sample has the same
initial history and opposite final bit with probability
\(2^{-r}/2=2^{-(r+1)}\). A union bound makes the probability of
at least one such conflicting sample at most \((w-1)/2^{r+1}\).
If there is no conflicting sample, every sample in that history group
has the same final bit, and the empirical optimum satisfies sample
\(j\). Therefore the expected success indicator of index \(j\)
is at least \(1-(w-1)/2^{r+1}\). Averaging over all indices
proves the expectation bound. No independence among these optimum-
success indicators is assumed.

Suppose the proposed accuracy event has probability at least
\(1-\delta\). On it \(\widehat v_*\leq1/2+\varepsilon\);
outside it \(\widehat v_*\leq1\). Hence
\[
 E[\widehat v_*]\leq
 (1-\delta)(1/2+\varepsilon)+\delta
 =1/2+\varepsilon+\delta(1/2-\varepsilon).
\]
Combining lower and upper bounds and rearranging gives
\(w-1\geq2^r(1-2\varepsilon)(1-\delta)\), as claimed.
For the numerical case the right side is
\(1024(4/5)(19/20)=778.24\); thus \(w\geq779.24\), and
the smallest possible integer sample size is 780. This is a necessary
condition obtained from the alleged guarantee, not a construction
proving 780 samples sufficient.
\end{proof}
\end{annexitem}

\begin{annexitem}{annex:p15.ssat-random-formula-bound}{thm:p15-ssat-random-formula-bound}{Independent clause generation and an exact first moment}
\begin{proof}
Fix a complete assignment. For any chosen set of \(k\) distinct
variables, independent fair signs make all \(k\) literals false
with probability \(2^{-k}\). Thus the assignment satisfies one
random clause with probability \(1-2^{-k}\). The \(m\)
clause draws are independent, so its probability of satisfying all
of them is \((1-2^{-k})^m\), even though duplicate clauses may
be drawn. Sum these indicators over the \(2^n\) assignments and
use finite linearity of expectation to obtain the displayed first
moment. For nonnegative \(Z\), the pointwise inequality
\(Z\geq2^t\mathbf1_{Z\geq2^t}\) gives Markov's bound by
taking expectations and dividing by positive \(2^t\).

For \(k=3\), the bound is
\(\exp[(n-t)\log2-m\log(8/7)]\). It tends to zero under
the stated divergence condition. At equality of the two exponents
it is exactly one and proves no asymptotic small-probability claim.

For a specified variable, the probability of belonging to the uniform
\(k\)-subset is \(k/n\), since each of the \(n\) variables
has the same inclusion probability and the subset has \(k\) members.
It is absent from all independent clause draws with probability
\((1-k/n)^m\). Summing its unused indicator over all variables
gives \(n(1-k/n)^m\). For \(m=0\), the empty product is one,
including \(k=n\); every variable is unused. This identifies the
sampling model responsible for the exact formula rather than using
a with-replacement approximation without qualification.
\end{proof}
\end{annexitem}
